\chapter{The classifier ladder}
\label{ch:classifiers}

The ladder is fixed in build order and was not reordered: a rule table, then an
n-gram logistic regression, then, only if the measured ceiling justified it, a
language model. The ordering is a claim about what a project owes its users, and
it is defended here rather than assumed.

\section{Why rules first}

A rule table is a table of regular expressions with a tier and a signal name on
every row. It has no probabilities and does not pretend to. It is also the
version of the tool that is genuinely useful to somebody who just wants their
checkout page fixed, and it shipped before any machine learning existed.

That ordering means the project can answer a question most projects cannot: is
the model solving a problem anybody has? A project that ships the model first
and the product last has no baseline to answer it with, and every subsequent
measurement is a comparison against nothing.

\subsection{Ties are the safety mechanism, not a shortcoming}

The rule engine has six precedence tiers. A tier decides when it yields a unique
match and otherwise falls through, and a tie surviving every tier is
\code{UNKNOWN}. That reads like a tidy up clause. It is the single most load
bearing rule in the engine.

German writes \emph{Stra\ss{}e und Hausnummer} for both a whole street address
and the first of several address lines. French writes \emph{Adresse} for both. A
field labelled only \emph{Password} is a login field on one page and a
registration field on the next. In each of those the wording genuinely does not
distinguish two labels, and the correct answer is that the tool does not know.

Once that is accepted the table changes shape. Instead of writing one pattern
per label and hoping they do not overlap, the overlaps become deliberate: both
labels carry the same phrase at the same weight, the tier ties by construction,
and the engine falls through to a tier that might separate them. The property
that a correct page produces no finding above the lowest severity fell out of
that design rather than being tuned into existence, and it passed the first time
it was run across the entire correct markup slice.

The alternative would have been to break each tie by picking a favourite. That
produces a confident instruction to write the wrong token into somebody's
production markup, and it would have looked like better coverage on every metric
except the one that matters.

\subsection{An honest dead rule}

The specification names the Japanese postal mark among the strings the postal
code label should match. It cannot fire: normalisation treats a lone symbol
character as a delimiter, so the mark never reaches a token stream, and the
Japanese postal field is reached through its word form instead.

The rule is in the table anyway, with a comment saying it cannot fire and why,
and a test asserting both halves of that. Making it reachable means changing what
counts as a token character, which changes every golden snapshot and every
committed fixture expectation. That is a deliberate decision with a cost rather
than a quiet fix, and it is recorded as one.

\section{Why an n-gram logistic regression before a transformer}

Three reasons, in the order they matter.

\textbf{It fits in a wheel.} The exported model plus its vocabulary and
calibration is small enough to install without comment, and inference needs an
array library and a runtime and nothing else. A tool that requires a large
download or a GPU to audit a form is a tool nobody installs.

\textbf{It is inspectable.} A linear model's weights name the character and word
n-grams behind a prediction, which is exactly what law 1 requires a finding to
carry. That property is not incidental to the design; it is why the rung exists
at that height.

\textbf{It sets a floor that a bigger model has to beat.} Without it, any
transformer result is a comparison against a rule table alone, and the
interesting question, whether the extra capacity buys anything on a task whose
whole input is a short list of attribute strings, never gets asked.

\subsection{Features}

Five blocks: character n-grams over the normalised token stream, word n-grams
over the same, one hot categorical signals, option shape booleans, and
structural buckets. Table~\ref{tab:training} gives the realised widths.

\begin{table}[htbp]
  \centering
  \caption{The training run that produced the shipped model.}
  \label{tab:training}
  \begingroup\small\setlength{\tabcolsep}{4pt}
\begin{tabular}{lr}
\toprule
\textbf{Training quantity} & \textbf{Value} \\
\midrule
Train forms & 500 \\  % results: models/dev_metrics.json
Train rows & 5402 \\  % results: models/dev_metrics.json
Dev forms & 120 \\  % results: models/dev_metrics.json
Dev rows & 1596 \\  % results: models/dev_metrics.json
Excluded forms & 100 \\  % results: models/dev_metrics.json
Classes fitted & 42 \\  % results: models/dev_metrics.json
Feature width & 7986 \\  % results: models/train_manifest.json
Character n-gram features & 5760 \\  % results: models/train_manifest.json
Word n-gram features & 2176 \\  % results: models/train_manifest.json
Categorical features & 32 \\  % results: models/train_manifest.json
Option shape features & 5 \\  % results: models/train_manifest.json
Structural features & 13 \\  % results: models/train_manifest.json
Chosen inverse regularisation & 8.0000 \\  % results: models/train_manifest.json
ONNX opset & 26 \\  % results: models/train_manifest.json
Dev macro-F1 & 0.7866 \\  % results: models/dev_metrics.json
Dev accuracy & 0.8421 \\  % results: models/dev_metrics.json
\bottomrule
\end{tabular}
\endgroup
\par\smallskip\noindent{\footnotesize Source: \rpath{models/dev_metrics.json}}
\par\noindent{\footnotesize Source: \rpath{models/train_manifest.json}}

\end{table}

\begin{table}[htbp]
  \centering
  \caption{The pre-registered hyperparameter sweep, on development macro-F1.
  The grid was registered before the model existed and was not widened
  afterwards.}
  \label{tab:sweep}
  \begingroup\small\setlength{\tabcolsep}{4pt}
\begin{tabular}{rr}
\toprule
\textbf{Inverse regularisation} & \textbf{Dev macro-F1} \\  % results: models/dev_metrics.json
\midrule
0.25 & 0.7308 \\  % results: models/dev_metrics.json
0.50 & 0.7375 \\  % results: models/dev_metrics.json
1.00 & 0.7505 \\  % results: models/dev_metrics.json
2.00 & 0.7612 \\  % results: models/dev_metrics.json
4.00 & 0.7660 \\  % results: models/dev_metrics.json
8.00 & 0.7690 \\  % results: models/dev_metrics.json
\midrule
chosen & 8.00 \\  % results: models/dev_metrics.json
\bottomrule
\end{tabular}
\endgroup
\par\smallskip\noindent{\footnotesize Source: \rpath{models/dev_metrics.json}}

\end{table}

\textbf{The sweep chose the edge of its own grid.} Table~\ref{tab:sweep} shows
the chosen value sitting at the largest value the grid offered. The grid was
pre-registered, so it was not widened afterwards: widening a pre-registered grid
because the answer landed at its edge is how a sweep becomes a search for a
number. The consequence is recorded instead, in the model card and here, and it
is that this model may be less regularised than a wider grid would have chosen.
A later phase that revisits it re-registers first.

\subsection{The export route, and why the fallback was the design}

The plan anticipated a fight with the converter for text vectorizers
\cite{skl2onnx} and named a fallback: move the whole text pipeline into hand
written code, export only the linear layer, and take that quickly rather than
spending a day coaxing a converter.

The fight never started, because the feature set was never convertible. Three of
the five feature blocks are categorical one hots, option shape booleans and
structural buckets, which no standard transformer produces. Putting the whole
featurisation in the graph would have meant writing custom converters for three
bespoke transformers \emph{and} trusting the vectorizer converters, which is a
larger version of exactly the risk the fallback exists to avoid.

So the fallback's hand written featuriser is not a retreat from the design. It
is the only design in which one implementation serves both training and
inference, which the feature specification required from the beginning. The
second argument arrived for free and will matter longer: a vectorizer anywhere in
the inference path puts a full machine learning framework on the dependency list
of a tool people install with one command \cite{scikit-learn,onnxruntime}.

The cost is real and is tested rather than waved past. A hand written analyser is
slower than a compiled one, and its faithfulness to the reference analyser is a
claim rather than an identity, so the claim is asserted character for character
over inputs including empty strings, words shorter than the n-gram length, and
Japanese text. The detail an independent reimplementation gets wrong is the
short word rule: a word shorter than the n-gram length yields its padded self
once rather than once per length, and getting that wrong silently triples the
weight of every two letter token.

\begin{table}[htbp]
  \centering
  \caption{The parity gate: the exported graph against the training time
  pipeline, on the full development split.}
  \label{tab:parity}
  \begingroup\small\setlength{\tabcolsep}{4pt}
\begin{tabular}{lr}
\toprule
\textbf{Parity quantity} & \textbf{Value} \\
\midrule
Dev rows compared & 1596 \\  % results: models/dev_metrics.json
Argmax agreements & 1596 \\  % results: models/dev_metrics.json
Agreement rate & 1.0000 \\  % results: models/dev_metrics.json
Largest absolute probability delta & 0.000000477 \\  % results: models/dev_metrics.json
Tolerance & 0.000010000 \\  % results: models/dev_metrics.json
Passed & yes \\
\bottomrule
\end{tabular}
\endgroup
\par\smallskip\noindent{\footnotesize Source: \rpath{models/dev_metrics.json}}

\end{table}

Table~\ref{tab:parity} is the gate that makes the export believable. The opset
was chosen by probing downward from the newest the installed toolchain defines
and taking the first that both converts and opens in a session; the converter
warns at every attempt that the requested opset is above the newest it has been
tested against. That warning is recorded rather than filtered, because it is
precisely the case where a successful export proves nothing and only the parity
gate proves something.

\subsection{A graph that cannot read its own weights}

Worth recording because it changed a design. The exported graph holds its
coefficients, intercepts and class order in a node's \emph{attributes} rather
than in graph initialisers, and the runtime does not hand a caller a node's
attributes. So the inference path cannot read its own weights, and naming the
n-grams behind a prediction would otherwise have required the full model format
library at runtime to produce one line of a report.

The resolution is a separate evidence table, written by the same training run,
holding each class's highest weighted features. A derived file is a drift risk,
so the parity suite reads the weights back out of the graph and asserts the
table agrees with them. A stale evidence file now fails the build rather than
naming the wrong n-grams in a finding, which would have satisfied law 1 in form
and violated it in substance.

\section{Calibration, and the part that got worse}

\begin{table}[htbp]
  \centering
  \caption{What calibration did to the model on the development split,
  including the summary statistic that moved the wrong way.}
  \label{tab:calibration}
  \begingroup\small\setlength{\tabcolsep}{4pt}
\begin{tabular}{lr}
\toprule
\textbf{Calibration quantity} & \textbf{Value} \\
\midrule
Classes fitted & 42 \\  % results: models/dev_metrics.json
Isotonic switchover, dev positives & 100 \\  % results: models/dev_metrics.json
Classes calibrated by identity & 2 \\  % results: models/dev_metrics.json
Classes calibrated by isotonic & 1 \\  % results: models/dev_metrics.json
Classes calibrated by sigmoid & 39 \\  % results: models/dev_metrics.json
Expected calibration error before & 0.0281 \\  % results: models/dev_metrics.json
Expected calibration error after & 0.0454 \\  % results: models/dev_metrics.json
Dev macro-F1 before calibration & 0.7690 \\  % results: models/dev_metrics.json
Dev macro-F1 after calibration & 0.7866 \\  % results: models/dev_metrics.json
\bottomrule
\end{tabular}
\endgroup
\par\smallskip\noindent{\footnotesize Source: \rpath{models/dev_metrics.json}}

\end{table}

\begin{figure}[htbp]
  \centering
  \includegraphics[width=.62\textwidth]{calibration_ngram.pdf}
  \caption[Reliability of the n-gram engine on the development split]{%
  Reliability of the n-gram engine on the development split, before and
  after calibration. Only this engine appears: the rule baseline has no
  probability at all, and the language model's confidence is self reported, so
  neither belongs on an axis of calibrated probability. Generated by
  \code{scripts/make\_report\_figures.py} from
  \rpath{models/dev_metrics.json}.}
  \label{fig:calibration}
\end{figure}

Expected calibration error was \emph{lower before calibration than after it}
(Table~\ref{tab:calibration}, Figure~\ref{fig:calibration}). The uncalibrated
model on this corpus is already close to honest, and refitting a one versus rest
sigmoid per class on a development split of one template per family, then
renormalising, moves classes that were already well behaved.

The calibrator stayed, for two reasons that are both stated rather than assumed.
Macro-F1 is higher with it. And the argument for calibrating at all is about the
number a single finding quotes to a developer, not about an average gap across a
split. What changed is the model card, which now says the error went the wrong
way and by how much. Dropping the calibrator because its summary statistic got
worse, without re-registering, would be tuning on the thing being measured.

One class reached the isotonic switchover, the label covering search boxes and
consent checkboxes. Every other class is calibrated by a sigmoid, and a handful
with no development positives at all are uncalibrated and named as such in the
card.

\section{The third rung, and what it is not}

The third engine is a twelve billion parameter instruction model at four bit
quantization, served locally through a runtime that enforces a JSON response
schema at decode time \cite{ollama-api}. It is optional at runtime, never
required for an audit, and never exercised by continuous integration.

\begin{table}[htbp]
  \centering
  \caption{What each engine says about itself. A run log records this block,
  because a comparison table that named neither the model tag nor the
  quantization would not say what it had measured.}
  \label{tab:engine-describe}
  \begingroup\small\setlength{\tabcolsep}{4pt}
\begin{tabular}{ll}
\toprule
\textbf{Engine property} & \textbf{Value} \\
\midrule
rules confidence kind & \texttt{tier} \\
ngram confidence kind & \texttt{calibrated-probability} \\
ngram execution provider & \texttt{CPUExecutionProvider} \\
ngram ONNX opset & \texttt{26} \\  % results: experiments/results/bench/2026-08-26T20-32-48Z_bench_6457ac7/benchmark.json
llm confidence kind & \texttt{self-reported} \\
llm model tag & \texttt{mistral-nemo:12b-instruct-2407-q4\_K\_M} \\  % results: experiments/results/bench/2026-08-26T20-32-48Z_bench_6457ac7/benchmark.json
llm quantization & \texttt{q4\_K\_M} \\  % results: experiments/results/bench/2026-08-26T20-32-48Z_bench_6457ac7/benchmark.json
llm temperature & \texttt{0.0} \\  % results: experiments/results/bench/2026-08-26T20-32-48Z_bench_6457ac7/benchmark.json
llm batch cap & \texttt{40} \\  % results: experiments/results/bench/2026-08-26T20-32-48Z_bench_6457ac7/benchmark.json
llm structured output & \texttt{json\_schema passed with every request} \\
llm prompt version & \texttt{1.0.0} \\  % results: experiments/results/bench/2026-08-26T20-32-48Z_bench_6457ac7/benchmark.json
\bottomrule
\end{tabular}
\endgroup
\par\smallskip\noindent{\footnotesize Source: \rpath{experiments/results/bench/2026-08-26T20-32-48Z_bench_6457ac7/benchmark.json}}

\end{table}

Three design points decide whether the comparison measures anything at all.

\textbf{The declaration is dropped from what the model sees.} Leaving it in lets
the model read the answer off the page on every clean tier form. The obvious
check, a substring search for the declared token in the serialised payload, is
wrong in both directions: a field named \code{email} that also declares
\code{autocomplete="email"} would trip it, and that combination is common enough
in a clean corpus that the check would have fired on the first real run and been
switched off, while a declaration reaching the payload through a derived value
would not trip it at all. So the property checked is \emph{independence}: pruning
a descriptor and pruning the same descriptor with its declaration stripped must
produce identical payloads. That has neither failure mode and is asserted on
every request rather than reviewed once.

\textbf{The keep list is a floor, not a ceiling.} The specification's list of
what pruning keeps omits several signals the featuriser feeds to the n-gram
model. Withholding a signal one engine has from the other tilts the comparison
exactly as far as handing over the whole document would, in the opposite
direction and with a better conscience. The rule applied is that the language
model sees what the featuriser sees and nothing more.

\textbf{Its threshold block is zero and zero, and that was registered before the
run.} The decision procedure is engine agnostic by confidence value and needs
two numbers, and the specification forbids a self reported confidence from
feeding the threshold policy. The only pair satisfying both is zero and zero.
The two consequences were written down in advance: the language model gets no
threshold protection and accuses on every committed prediction, and the low
confidence finding becomes unreachable for it. Deriving a boundary on the self
reported scale against a precision target, the way the n-gram engine's was
derived, would probably have raised its finding level precision. That it would
have flattered the engine is exactly why it had to be refused before the
measurement rather than after.
